\section{Reproducibility and fairness}
\label{sec:platform-repro}

The pre-study (Section~\ref{sec:prestudy-repro}) asked for fixed seeds and
repeated runs. In a benchmark that compares methods, fixing a seed is
necessary but not sufficient: the comparison is fair only if every method is
trained and evaluated on the same problem instances from the same starting
points, and it is reproducible only if the code that produced a number can be
identified afterwards. The platform provides both as mechanisms rather than
as conventions the user must remember.

\paragraph{One seed, threaded everywhere.}
A config carries one top-level \texttt{seed}. The runner copies it into the
evaluation flags, and every adapter reads it for training; all six
libraries' \texttt{train.py} and \texttt{evaluate.py} accept \texttt{--seed},
added wherever a script did not already expose one. On the training side the seed fixes weight initialisation and the
sampling of unroll segments; for the model-free optimizers this was verified
to be bit-exact across separate processes. On the evaluation side it fixes which
instances are evaluated and where each trajectory starts.

\paragraph{Every method sees the same instances.}
For the LASSO experiments, evaluation never draws from a mutable random
state. The model-free evaluation scripts walk the test split with a plain
integer cursor, batch by batch, covering every instance exactly once; the
model-based test path applies no shuffling. Consequently the number of random
draws a method consumed during training, which differs between methods,
cannot change which instances it is scored on. This is not left to code
reading: every LASSO-capable adapter reports an \emph{instances fingerprint},
a hash of the $b$, $x_\mathrm{true}$ and $x_0$ arrays its evaluation script
actually consumed, in the order it consumed them, which is stored with the
run's metrics and therefore in every sweep's result rows. A test suite
(\texttt{Benchmarking/tests/}) asserts that all methods in a cell report
the same fingerprint and that the count equals the size of the test split,
both live on a small generated fixture and post hoc on any completed sweep;
for the sweeps of Chapter~\ref{ch:exp2} every cell agrees. The fingerprint
is problem-agnostic (it hashes named arrays; what they are is the adapter's
choice), so a new problem is covered by adding one entry to the test matrix.

\paragraph{The same run twice.}
The same test suite evaluates every method twice with the same seed and
requires identical trajectories (bit-exact for the file-based problems), and,
for problems whose instances are drawn from the seed, that two seeds do not
produce the same run.

\paragraph{Every method starts from the same point.}
Each split file has a seeded sibling \texttt{*\_x0.npy} holding one starting
vector $x_0$ per instance (Section~\ref{sec:platform-data}). Originally the
model-free optimizers drew a random $x_0$ at evaluation time and the unrolled
networks implicitly started from $x_0 = 0$, so early iterates were not
comparable across methods. Now all four Experiment-2 methods consume the same
$x_0$: the model-free optimizee is initialised from the file, and for LISTA
and ALISTA the vector is appended to the network input and consumed by layer
zero. For LISTA this required a fixed, non-trainable first step
$x_1 = S_{\theta_1}\big((I - BA)\,x_0 + B b\big)$, i.e.\ one ISTA step from
$x_0$, so that the network keeps exactly $T$ learnable layers; ALISTA's layer
formula already generalises to a nonzero start. The layer-zero outputs were
checked against an independent hand computation of one ISTA step, and the
model-free optimizers' first recorded objective value against the mean
$f(x_0)$ computed directly from the file.

\paragraph{Provenance.}
Every run records the exact subprocess command lines, wall-clock timings and
the seed alongside its metrics. A sweep additionally records the git commit,
branch and working-tree status, a diff of uncommitted changes, and a SHA-256
hash of every untracked Python file, together with a snapshot of the sweep's
own configuration. The reason for going beyond a commit hash is practical:
during a multi-day sweep the code that produced a row is usually not yet
committed, and a commit hash alone would identify the wrong code.

\paragraph{Environment.}
All experiments reported in this thesis ran inside one container image built
from the repository's \texttt{Dockerfile}: Python 3.12 with TensorFlow 2.21
and Keras 3.15 in an isolated virtual environment (the release's single
PyTorch library gets a second one, because the two frameworks pin conflicting
CUDA packages), on a single 8\,GB consumer GPU. A run is resumable: the sweep
driver skips cells whose rows already exist, so a stopped container continues
where it left off with no change in results.

\paragraph{A known gap.}
Seeding fixes \emph{which} problem is evaluated, not every floating-point
detail of evaluating it. For neural-network optimizees in the L2O-Scale family
(MNIST classifiers) two runs with the same seed differ by about $0.1\%$,
attributable to non-deterministic kernels and Keras layer initialisation that
bypasses the seed argument; this predates the platform. The LASSO problems used
in this thesis are unaffected: their instances are read from files and their
evaluation is a deterministic pass.
